\begin{abstract}
Lung nodule detectors rank candidate locations in chest CT more accurately when they examine a wider region around each candidate, but it is not known why: the wider view could show the surrounding anatomy, such as vessels, or simply more of the nodule. We separate these explanations with a controlled experiment. Keeping the candidate locations, the network and its training fixed, we change only which parts of the original scan the network may see: a small sphere around the candidate, the sphere plus a thin surrounding shell, or the full view. On a standard benchmark, restoring the thin shell alone recovers about 70\% of the benefit of the full view, raising mean sensitivity from 83.6\% to 87.6\%. This benefit requires dense coverage of the shell rather than a few well-chosen samples of it, and it concentrates on large or off-centre nodules whose tissue the shell reaches. Much of the value of wider context therefore comes from seeing more of the lesion itself, which suggests choosing the field of view from the size of the target lesions rather than from the scale of the surrounding anatomy.
\end{abstract}
